A mi directora de tesis, la Dra. Noemí Patricia Kisbye, por su inmensa paciencia, su guía experta y su dedicación a lo largo de todo este proyecto. Gracias por impulsarme a mantener el rigor matemático y analítico, y por brindarme las herramientas necesarias para llevar a buen puerto esta investigación.

A la Facultad de Matemática, Astronomía, Física y Computación (FAMAF) y a la Universidad Nacional de Córdoba (UNC), por haberme brindado una educación pública, gratuita y de excelencia. A todos los profesores que, a lo largo de la carrera, me transmitieron su pasión por las Ciencias de la Computación y formaron mi criterio profesional.

A mis padres, Sandra Delia Fiore y Hector Omar Bordón, y a mi hermano, Martín Ezequiel Bordón, por su apoyo incondicional, por el inmenso esfuerzo que han hecho para que yo pueda llegar hasta acá y por ser mi pilar fundamental durante todos estos años. Este logro es tan mío como de ustedes.

A mis compañeros de cursada, y muy especialmente a todo Tiesos FC, porque aunque son muchos para nombrarlos a cada uno, las incontables horas compartidas entre apuntes, algoritmos y mates hicieron que el camino universitario haya sido una experiencia inolvidable.

A mis amigos, todos aquellos que conocí en circunstancias tan peculiares y distintas pero especialmente a Nueva Chicago y todo El Descansadero (he aquí una mención a Juan Cruz Roldán y Nazareno Kohan Boc con quienes nos comprometimos en hacernos aparecer en algún escrito del otro), por cada baile, juego, ecuación, código, risa o momento compartido que me recuerdan que siempre hay algo más por aprender.

Un agradecimiento muy especial (y lleno de cariño) a Katherina Bordón e Inari Bordón. Gracias por hacerme compañía en las largas madrugadas frente a la pantalla, por su amor incondicional y por ser mi cable a tierra durante todo este proceso.

Finalmente, a todos aquellos que de una u otra forma contribuyeron a que este trabajo sea posible. ¡Muchas gracias!
